\section{K$_3$C$_{60}$ details}

The K$_3$C$_{60}$ case study in the main text combines an archived near-Fermi band-context figure with a bundled mode-resolved curvature export. The archived band plot is \texttt{figures/supplementary/K3C60-bands.pdf}. The reproducible curvature diagnostic uses \texttt{data/k3c60/tmp.agr}, whose embedded axis metadata identify the plotted quantity as $\partial^2 E_n(k)/\partial q^2$ along the $\Gamma$--X--L--W--$\Gamma$ path. The archived companion figure and caption identify the representative phonon mode as \texttt{m120} at $1090.35~\mathrm{cm}^{-1}$.

The curvature-window analysis shown in the main-text K$_3$C$_{60}$ figure and in Fig.~\ref{fig:s5_threshold_bootstrap} is generated by \texttt{scripts/make\_fig2\_k3c60\_edt.py}. The nominal analysis uses Bands 208 and 209 from the archived numbering, a momentum slice corresponding to indices 150--281 of the exported trace, a local quadratic fit half-window of 18 points, and a residual threshold $R_{\rm th}=0.08$. These choices are not hidden analysis knobs: each of them is varied in the robustness scans reported in Fig.~\ref{fig:s5_threshold_bootstrap}.

Because the bundled archive contains one exported curvature dataset rather than the full family of displaced calculations, the reproducible workflow starts from the already-computed curvature traces and does not reconstruct them from multiple displacement amplitudes. The manuscript therefore treats displacement-amplitude robustness conservatively: it is explicitly marked as unavailable in the archived reproducibility bundle rather than inferred from unarchived calculations.
